\clearpage
\SourceLeaf{610}
\ArithmeticHead{581}{R}{INHALT.}
\begingroup
\newcommand{\ContentsRow}[5]{\par\noindent\hangindent=#1\hspace*{#1}#2\hspace{.4em}\mbox{\footnotesize #3}\hfill\mbox{#4\quad #5}\par}
\noindent\textls[120]{Zur Theorie des arithmetisch-geometrischen Mittels.}\par
\ContentsRow{1em}{I.\quad Specimen termini medii si nomine uti licet arithmetico-geometrici}{.\,.\,.\,.\,.\,.\,.\,.}{Seite}{172}
\ContentsRow{1em}{II.\quad Reihenentwicklungen und Beziehungen zum Ellipsenumfang}{.\,.\,.\,.\,.\,.\,.\,.\,.\,.}{—}{177}
\ContentsRow{1em}{III.\quad Differential- und Funktionalgleichungen}{.\,.\,.\,.\,.\,.\,.\,.\,.\,.\,.\,.\,.\,.\,.\,.\,.}{—}{181}
\ContentsRow{1em}{IV.\quad Investigatio functionum quae ex evolutione mediorum arithmetico-geometricorum oriuntur}{}{—}{184}
\ContentsRow{1em}{V.\quad Theoria sinus lemniscatici universalissime accepti}{.\,.\,.\,.\,.\,.\,.\,.\,.\,.\,.\,.\,.\,.}{—}{194}
\ContentsRow{1em}{VI.\quad Differentiatio mediorum arithmetico-geometricorum}{.\,.\,.\,.\,.\,.\,.\,.\,.\,.\,.\,.\,.}{—}{207}
\ContentsRow{1em}{VII.\quad Algorithmi ad media arithmetico-geometrica pertinentes}{.\,.\,.\,.\,.\,.\,.\,.\,.\,.\,.}{—}{209}
\ContentsRow{1em}{VIII.\quad Der bilineare Algorithmus}{.\,.\,.\,.\,.\,.\,.\,.\,.\,.\,.\,.\,.\,.\,.\,.\,.\,.\,.\,.\,.\,.}{—}{213}
\ContentsRow{1em}{IX.\quad Die Beziehung zwischen den unendlich vielen Werten des\linebreak arithmetisch-geometrischen Mittels. Unmittelbare Anwendung der Theorie auf die elliptischen Transzendenten und auf die Rektifikation der Ellipse}{.\,.\,.\,.\,.\,.\,.\,.\,.\,.\,.\,.\,.\,.\,.\,.}{—}{217}
\par\medskip\noindent\hspace*{1em}\emph{Briefwechsel:}\par
\ContentsRow{2em}{1.\ \textsc{Pfaff} an \textsc{Gauss}, 1799 November 24}{.\,.\,.\,.\,.\,.\,.\,.\,.\,.\,.\,.\,.\,.\,.\,.}{—}{232}
\ContentsRow{2em}{2.\ \textsc{Pfaff} an \textsc{Gauss}, 1800 Dezember 8}{.\,.\,.\,.\,.\,.\,.\,.\,.\,.\,.\,.\,.\,.\,.\,.\,.}{—}{234}
\ContentsRow{3em}{Anhang zu 2: aus \textsc{Schumachers} \emph{Gaussiana}}{.\,.\,.\,.\,.\,.\,.\,.\,.\,.\,.\,.\,.\,.}{—}{236}
\ContentsRow{2em}{3.\ \textsc{Gauss} an \textsc{Bessel}, 1805 September 3}{.\,.\,.\,.\,.\,.\,.\,.\,.\,.\,.\,.\,.\,.\,.\,.}{—}{237}
\ContentsRow{2em}{4.\ \textsc{Schumacher} an \textsc{Gauss}, 1808 April 2}{.\,.\,.\,.\,.\,.\,.\,.\,.\,.\,.\,.\,.\,.\,.\,.}{—}{242}
\ContentsRow{2em}{5.\ \textsc{Gauss} an \textsc{Schumacher}, 1808 September 17}{.\,.\,.\,.\,.\,.\,.\,.\,.\,.\,.\,.\,.\,.}{—}{243}
\ContentsRow{2em}{6.\ \textsc{Schumacher} an \textsc{Gauss}, 1816 April 5}{.\,.\,.\,.\,.\,.\,.\,.\,.\,.\,.\,.\,.\,.\,.\,.}{—}{245}
\ContentsRow{2em}{7.\ \textsc{Gauss} an \textsc{Schumacher}, 1816 April}{.\,.\,.\,.\,.\,.\,.\,.\,.\,.\,.\,.\,.\,.\,.\,.\,.}{—}{247}
\ContentsRow{2em}{8.\ \textsc{Gauss} an \textsc{Bessel}, 1828 März 30}{.\,.\,.\,.\,.\,.\,.\,.\,.\,.\,.\,.\,.\,.\,.\,.\,.\,.}{—}{248}
\ContentsRow{3em}{Anhang zu 8: a.\ \textsc{Pfaff} an \textsc{Gauss}, 1824 Oktober 20}{.\,.\,.\,.\,.\,.\,.\,.\,.\,.\,.}{—}{249}
\ContentsRow{5em}{b.\ \textsc{Gauss} an \textsc{Pfaff}, 1825 März 21}{.\,.\,.\,.\,.\,.\,.\,.\,.\,.\,.\,.}{—}{250}
\noindent\hspace*{1em}Bemerkungen.\par
\ContentsRow{2em}{Abriß der Theorie des arithmetisch-geometrischen Mittels}{.\,.\,.\,.\,.\,.\,.\,.\,.\,.}{—}{251}
\ContentsRow{2em}{Erläuterungen zu I}{.\,.\,.\,.\,.\,.\,.\,.\,.\,.\,.\,.\,.\,.\,.\,.\,.\,.\,.\,.\,.\,.\,.\,.}{—}{257}
\ContentsRow{2em}{Erläuterungen zu II}{.\,.\,.\,.\,.\,.\,.\,.\,.\,.\,.\,.\,.\,.\,.\,.\,.\,.\,.\,.\,.\,.\,.\,.}{—}{261}
\ContentsRow{2em}{Erläuterungen zu III}{.\,.\,.\,.\,.\,.\,.\,.\,.\,.\,.\,.\,.\,.\,.\,.\,.\,.\,.\,.\,.\,.\,.}{—}{266}
\ContentsRow{2em}{Erläuterungen zu IV}{.\,.\,.\,.\,.\,.\,.\,.\,.\,.\,.\,.\,.\,.\,.\,.\,.\,.\,.\,.\,.\,.\,.}{—}{268}
\ContentsRow{2em}{Erläuterungen zu V}{.\,.\,.\,.\,.\,.\,.\,.\,.\,.\,.\,.\,.\,.\,.\,.\,.\,.\,.\,.\,.\,.\,.\,.}{—}{274}
\ContentsRow{2em}{Erläuterungen zu VI und VII}{.\,.\,.\,.\,.\,.\,.\,.\,.\,.\,.\,.\,.\,.\,.\,.\,.\,.\,.\,.}{—}{278}
\ContentsRow{2em}{Erläuterungen zu VIII}{.\,.\,.\,.\,.\,.\,.\,.\,.\,.\,.\,.\,.\,.\,.\,.\,.\,.\,.\,.\,.\,.\,.}{—}{279}
\ContentsRow{2em}{Erläuterungen zu IX}{.\,.\,.\,.\,.\,.\,.\,.\,.\,.\,.\,,\,.\,.\,.\,.\,.\,.\,.\,.\,.\,.\,.}{—}{280}
\ContentsRow{2em}{Erläuterungen zum Briefwechsel}{.\,.\,.\,.\,.\,.\,.\,.\,.\,.\,.\,.\,.\,.\,.\,.\,.\,.\,.}{—}{283}
\par\medskip\noindent\textls[120]{Zur Theorie der transscendenten Functionen gehörig.}\par
\ContentsRow{1em}{I.\quad Entwickelungen und Umformungen der unendlichen Reihen und Produkte}{.\,.\,.\,.\,.}{—}{287}
\ContentsRow{1em}{II.\quad Lösung des Umkehrproblems für das elliptische Integral erster Gattung}{.\,.\,.\,.\,.\,.}{—}{308}
\ContentsRow{1em}{III.\quad In den kleinsten Teilen ähnliche Abbildung der Fläche einer Ellipse auf die Fläche eines Kreises}{.\,.\,.\,.\,.\,.\,.\,.\,.\,.\,.\,.\,.\,.\,.\,.\,.\,.\,.\,.\,.\,.\,.\,.\,.}{—}{311}
\noindent\hspace*{1em}Bemerkungen.\par
\ContentsRow{2em}{Erläuterungen zu I und II}{.\,.\,.\,.\,.\,.\,.\,.\,.\,.\,.\,.\,.\,.\,.\,.\,.\,.\,.\,.\,.\,.}{—}{320}
\ContentsRow{2em}{Erläuterungen zu III}{.\,.\,.\,.\,.\,.\,.\,.\,.\,.\,.\,.\,.\,.\,.\,.\,.\,.\,.\,.\,.\,.\,.\,.}{—}{323}
\endgroup
